% Training methods for r21vgg2k_vgg16_bn_aa5_s42.
% New citation keys zhang2019making and cao2018vggface2 should be matched
% to the manuscript bibliography. Other keys follow the supplied old text.

A single VGG16 network with batch normalization was trained from scratch jointly on face, house, and object classes using a unified softmax classification head. The network received $180\times180\times3$ inputs produced by foveation and log-polar transformation. The convolutional features were globally averaged to produce a 512-dimensional vector, followed by a fully connected layer of 256 sigmoid units with temperature $T=2.0$. These units provided deterministic continuous activations during training. Dropout of 0.3 was applied before the final classification layer, which contained 2,250 outputs. Training minimized cross-entropy loss using AdamW with an initial learning rate of $10^{-3}$, weight decay of $5\times10^{-2}$ (applied only to parameters with dimensionality greater than one), and a batch size of 256 fixation crops. CUDA training used bfloat16 mixed precision and channels-last memory format.

The network used antialiased downsampling following \cite{zhang2019making}. At each of the five pooling stages, $2\times2$ max pooling with stride one was followed by a fixed $5\times5$ binomial filter and downsampling by a factor of two. The filter was the normalized outer product of $[1,4,6,4,1]$ with itself, applied independently to each feature channel with reflection padding. This bin5 operation low-pass filtered the feature maps before reducing their spatial resolution.

Training examples were sampled with replacement using weights determined by the currently active curriculum stage. Faces, objects, building identities, and the generic house class were treated as four sampling groups. Each group's sampling probability was proportional to the square root of its number of active classes; for the generic house group, which contained a single class, the number of active training images was used instead. Within each group, class probabilities were proportional to the square root of the number of training images in that class, and fixation crops within a class were sampled uniformly. Each epoch consisted of 2,000 optimization steps, corresponding to 512,000 sampled fixation crops.

We used a curriculum-learning strategy, where the number of classes the model had to recognize was progressively increased during training (\cite{bengio2009curriculum}). Face identities were introduced in ten nested stages of 4, 8, 16, 32, 64, 128, 256, 512, 1,024, and 2,048 identities. The corresponding numbers of object classes were 0, 0, 4, 8, 16, 32, 64, 64, 64, and 64, and the numbers of building identities were 0, 0, 0, 4, 8, 16, 32, 64, 137, and 137. The generic house class was introduced at stage five, with 4, 8, 16, 32, 64, and 128 training images available in successive stages. Training lasted 2, 2, 4, 4, 8, 8, 16, 16, 32, and 32 epochs per stage, respectively, for a total of 124 epochs. Previously introduced classes remained available in subsequent stages, and classification was restricted to the active classes. A cosine learning-rate schedule was applied across the full curriculum, with a 500-step linear warm-up after each stage transition. Model initialization and training sampling used seed 42; object and building class orderings used seed 0, while face identities followed a fixed ordering described below.

Saliency fixation locations for each full image were precomputed, but the corresponding transformations were applied on-the-fly during training. Training augmentation consisted of random resized crops and independent random rotations between $-15^\circ$ and $15^\circ$, applied before foveation and log-polar transformation. Validation images were presented upright. After each epoch, logits were summed across the first 16 fixation views of each validation image and the highest-scoring active class was selected. The checkpoint with the highest validation accuracy within the final curriculum stage was retained, together with the final training checkpoint. Early stopping was configured after 20 epochs without validation improvement within a stage, with advancement to the next stage when applicable. The run completed all 124 epochs; the highest validation accuracy occurred at epoch 122.

\subsection{Data}
We trained a single network jointly on three categories: faces, objects, and houses in one unified label space. The house category included both individual building identities and a generic house class.

\textbf{Faces Training: }2,048 identities were selected from the training partition of VGGFace2 (\cite{cao2018vggface2}), with 409,581 training, 24,576 validation, and 24,576 test images. Each identity contributed 200 training images, except for one identity with 181; all identities contributed 12 validation and 12 test images. The label was the identity, so the face task was fine-level discrimination. Identities followed a fixed nested ordering: the first eight were manually selected, and the first 512 were assigned to the Caucasian group. The remaining identities were interleaved to produce a final set of 1,500 Caucasian and 548 other-group identities. These were operational group assignments based on FairFace predictions, combined into Caucasian, African, Asian, and Indian groups; the initial eight identities were manually labeled. Identities from the other groups first entered at stage nine.

\textbf{Objects: }64 ImageNet categories, with 65,689 training, 8,208 validation, and 8,218 test images, split approximately 80/10/10. The labels were object categories, so the object task was categorization rather than identification.

\textbf{Houses: }137 building identities from the ZuBuD database were used for training. Each building contributed four exterior views for training and one for validation, giving 548 training and 137 validation images. The label was the building identity. An additional generic house class used frontal photographs from the Houses-dataset. Its training pool contained 435 images, of which a nested subset of up to 128 was used according to the curriculum. A separate set of 100 images was used for validation. Together, these sources provided 137 building identity classes and one generic house class.

\subsection{Transformations}
Here we describe the details of the transformations applied to the images before presenting them to the network. Images were converted to RGB, resized so that the shorter side was 224 pixels, and center-cropped to $224\times224$ pixels. Fixation locations were computed from these images, and the network inputs were constructed around the sampled fixations.

\subsubsection{Salience Map}
A Gaussian center prior ($\sigma=W/4$, raised to the second power) was applied to each grayscale image to place greater weight on its center. The weighted image was log-polar transformed and standardized before applying a Gabor-based salience operator (\cite{yamada1995model}). The image was passed through a bank of $31\times31$ Gabor filters at two scales ($\lambda=4,8$, with $\sigma=0.5\lambda$), four orientations ($\theta=0,\pi/4,\pi/2,3\pi/4$), and two phases ($\psi=0,\pi/2$), with aspect ratio 0.5. Each quadrature phase pair was combined into a magnitude, giving eight magnitude maps. Their variance at each pixel defined the saliency map, assigning greater weight to regions with varied local structure. A ten-pixel boundary was excluded from the filtered log-polar maps.

Thirty-two fixation locations were sampled without replacement from the log-polar saliency map, with sampling probabilities proportional to saliency, and mapped back to Cartesian image coordinates. Sampling was restricted to the central 70\% of the image along each axis. Training used 16 fixation slots per image; each time a slot was sampled, its fixation was randomly selected from the corresponding pair among the 32 stored locations ($j$ or $j+16$, for $j=0,\ldots,15$). Validation used the first 16 stored locations.

\subsubsection{Cropping and Rotation}
Each image was cropped from $224\times224$ pixels to $180\times180$ pixels around the selected fixation, with zero padding where the crop extended beyond the image boundary. During training, a randomly positioned square covering between 60\% and 100\% of the fixation crop's area was resampled to $180\times180$ pixels using bilinear interpolation. Each crop was then independently rotated by an angle sampled uniformly between $-15^\circ$ and $15^\circ$. No random cropping or rotation was applied during validation.

\subsubsection{Foveation}
We simulated retinal foveation using a six-level Gaussian pyramid centered on the crop (\cite{Jiang2015}). Successive levels were blurred and downsampled, then resized to the original crop dimensions and blended according to distance from the center. The resulting image had its highest resolution at the center, with resolution decreasing toward the periphery.

\subsubsection{Log-Polarization}
Each foveated crop was transformed from Cartesian coordinates $(x,y)$ to log-polar coordinates $(\log r,\theta)$ around the crop center. Radial distance was sampled on a logarithmic scale and polar angle was sampled linearly over a full revolution, producing a $180\times180$ representation with angle along the rows and log radius along the columns. Sampling coordinates outside the crop were clamped to its boundary. The same transformation was applied to each color channel.
